\clearpage
\SourceLeaf{321}
\section*{BEMERKUNGEN.}
\pdfbookmark[1]{Bemerkungen zur geometrischen Seite der ternären Formen}{ternary-bemerkungen}
\BookRule{.15\textwidth}
Neben den vorstehenden Notizen, die die in der Besprechung von \textsc{Seeber}s Untersuchungen der ternären Formen dargestellten Gesichtspunkte teilweise weiterentwickeln, sind in der Handschrift mehrere eigene Kristallmessungen mit einem achtzölligen \textsc{Reichenbach}schen Theodoliten aufgezeichnet. Die einzelnen Protokolle enthalten jeweils das Beobachtungsdatum, aus dem hervorgeht, dass diese Untersuchung dem Juli $1831$ angehört.

Aus der Theorie der indefiniten ternären quadratischen Formen findet sich im handschriftlichen Nachlass nur der folgende, wahrscheinlich während der Ausarbeitung der Disqu. Arr. aufgezeichnete Satz
\begin{quote}
„Alle Transformationen der ternären Form
\[
\left(\begin{smallmatrix}1,&1,&-1\\0,&0,&0\end{smallmatrix}\right)
\]
in sich selbst werden durch die Formel
\[
\begin{array}{ccc}
\alpha\delta+\beta\gamma&\alpha\beta-\gamma\delta&\alpha\beta+\gamma\delta\\[.5ex]
\alpha\gamma-\beta\delta&\frac12(\alpha\alpha+\delta\delta-\beta\beta-\gamma\gamma)&\frac12(\alpha\alpha+\gamma\gamma-\beta\beta-\delta\delta)\\[.5ex]
\alpha\gamma+\beta\delta&\frac12(\alpha\alpha+\beta\beta-\gamma\gamma-\delta\delta)&\frac12(\alpha\alpha+\beta\beta+\gamma\gamma+\delta\delta)
\end{array}
\]
dargestellt, wobei $\alpha,\beta,\gamma,\delta$ so gewählt werden, dass $\alpha\delta-\beta\gamma=1$ gilt.“
\end{quote}

Es entstehen nämlich alle Transformationen, in denen die neun Koeffizienten ganze Zahlen sind, wenn für $\alpha,\beta,\gamma,\delta$ sowohl alle die der Bedingungsgleichung genügenden ganzen Zahlen und zwar zwei gerade und zwei ungerade gesetzt werden, als auch alle die ungeraden Vielfachen von $\surd\frac12$, die dieselbe Bedingungsgleichung $\alpha\delta-\beta\gamma=1$ erfüllen.

\SourceLeaf{322}
Zu Seite $309$. Chaux carbonatée équiaxe, inverse, contrastante und mixte sind die von \textsc{Hauy} (Traité de Minéralogie $1801$ Band II, S. $132$, $137$) gebrauchten Benennungen.

Die Tafel der Transformationen der Form $\left(\begin{smallmatrix}5,&5,&5\\1,&1,&1\end{smallmatrix}\right)$ enthält in der ersten vertikalen Spalte die Koeffizienten der Substitution, in der zweiten die dadurch entstandene neue Form, in der dritten die der letzteren Form entsprechende primitive, wenn diese nicht bereits selbst primitiv ist, und in der vierten deren adjungierte Form.
\begin{flushright}\textsc{Schering.}\end{flushright}
\BookRule{.30\textwidth}

